
\documentclass{article}
\usepackage{colortbl}
\usepackage{makecell}
\usepackage{multirow}
\usepackage{supertabular}

\begin{document}

\newcounter{utterance}

\centering \large Interaction Transcript for game `eqbench', experiment `eq\_bench', episode 12 with reflexion{-}Qwen3.5{-}27B.
\vspace{24pt}

{ \footnotesize  \setcounter{utterance}{1}
\setlength{\tabcolsep}{0pt}
\begin{supertabular}{c@{$\;$}|p{.15\linewidth}@{}p{.15\linewidth}p{.15\linewidth}p{.15\linewidth}p{.15\linewidth}p{.15\linewidth}}
    \# & \multicolumn{2}{c}{Player} && \multicolumn{2}{c}{Game Master} \\
    \hline

    \theutterance \stepcounter{utterance}  
    & & & \multicolumn{4}{p{0.6\linewidth}}{
        \cellcolor[rgb]{0.9,0.9,0.9}{
            \makecell[{{p{\linewidth}}}]{
                \texttt{\tiny{[P1$\langle$GM]}}
                \texttt{You have attempted this same task before and did not succeed. These are the plans you wrote after those attempts.} \\
\texttt{Trial \#1: New plan: My previous attempt failed because I assigned a non{-}zero score to "Fear" (2), which was not the most accurate interpretation of Lorelei's emotional state given the context of a heated political debate where she is actively arguing her points rather than expressing terror. While "Frustration" was correctly identified as the primary emotion, the inclusion of "Fear" diluted the accuracy of the prediction. Next time, I will strictly analyze the dialogue for the *dominant* emotional tone and assign a score of 0 to any emotion that is not clearly supported by the text, even if it seems remotely plausible. Specifically, I will recognize that Lorelei's tone is one of passionate disagreement and exasperation, not fear, and ensure only "Frustration" receives a high score while "Apathy," "Fear," and "Agreement" are all scored as 0.} \\
\\ 
\texttt{Trial \#2: New plan: My previous attempt failed because I assigned a non{-}zero score to "Fear" (2), which was not the most accurate interpretation of Lorelei's emotional state given the context of a heated political debate where she is actively arguing her points rather than expressing terror. While "Frustration" was correctly identified as the primary emotion, the inclusion of "Fear" diluted the accuracy of the prediction. Next time, I will strictly analyze the dialogue for the *dominant* emotional tone and assign a score of 0 to any emotion that is not clearly supported by the text, even if it seems remotely plausible. Specifically, I will recognize that Lorelei's tone is one of passionate disagreement and exasperation, not fear, and ensure only "Frustration" receives a high score while "Apathy," "Fear," and "Agreement" are all scored as 0.} \\
\\ 
\texttt{Trial \#3: New plan: My previous attempt failed because I assigned a non{-}zero score to "Fear" (2), which was not the most accurate interpretation of Lorelei's emotional state given the context of a heated political debate where she is actively arguing her points rather than expressing terror. While "Frustration" was correctly identified as the primary emotion, the inclusion of "Fear" diluted the accuracy of the prediction. Next time, I will strictly analyze the dialogue for the *dominant* emotional tone and assign a score of 0 to any emotion that is not clearly supported by the text, even if it seems remotely plausible. Specifically, I will recognize that Lorelei's tone is one of passionate disagreement and exasperation, not fear, and ensure only "Frustration" receives a high score while "Apathy," "Fear," and "Agreement" are all scored as 0.} \\
\\ 
\\ 
\\ 
\texttt{Your task is to predict the likely emotional responses of a character in this dialogue:} \\
\\ 
\texttt{Harold: I don't understand why you're so worked up about this election. It's not like the world is going to end if your candidate doesn't win.} \\
\texttt{Lorelei: That's not the point, Harold. It's about what this election represents. It's about the future of our planet, our society.} \\
\texttt{Harold: And you think your candidate has all the answers? } \\
\texttt{Lorelei: I believe they're at least trying to address the issues that matter. Climate change, social justice, healthcare...} \\
\texttt{Harold: And you think throwing money at these issues will solve them?} \\
\texttt{Lorelei: It's not about throwing money, Harold. It's about investing in solutions, in our future.} \\
\texttt{Harold: And who's going to pay for all these grand plans? Us, the taxpayers. } \\
\texttt{Lorelei: Isn't that the point of taxes? To fund public services and infrastructure?} \\
\texttt{Harold: Yes, but there's a limit. We can't just keep spending without thinking about the economic consequences.} \\
\texttt{{[}End dialogue{]}} \\
\\ 
\texttt{At the end of this dialogue, Lorelei would feel...} \\
\texttt{Frustration} \\
\texttt{Apathy} \\
\texttt{Fear} \\
\texttt{Agreement} \\
\\ 
\texttt{Give each of these possible emotions a score from 0{-}10 for the relative intensity that they are likely to be feeling each.} \\
\\ 
\texttt{You must output in the following format, including headings (of course, you should give your own scores), with no additional commentary:} \\
\\ 
\texttt{Frustration: <score>} \\
\texttt{Apathy: <score>} \\
\texttt{Fear: <score>} \\
\texttt{Agreement: <score>} \\
\\ 
\\ 
\texttt{{[}End of answer{]}} \\
\\ 
\texttt{Remember: zero is a valid score, meaning they are likely not feeling that emotion. You must score at least one emotion > 0.} \\
\\ 
\texttt{Your answer:} \\
            }
        }
    }
     \\ \\

    \theutterance \stepcounter{utterance}  
    & \multicolumn{4}{p{0.6\linewidth}}{
        \cellcolor[rgb]{0.561,0.737,0.561}{
            \makecell[{{p{\linewidth}}}]{
                \texttt{\tiny{[P1$\rangle$GM]}}
                \texttt{Frustration: 9} \\
\texttt{Apathy: 0} \\
\texttt{Fear: 0} \\
\texttt{Agreement: 0} \\
            }
        }
    }
    & & \\ \\

    \theutterance \stepcounter{utterance}  
    & & & \multicolumn{2}{p{0.3\linewidth}}{
        \cellcolor[rgb]{0.9,0.9,0.9}{
            \makecell[{{p{\linewidth}}}]{
                \texttt{\tiny{[GM$|$GM]}}
                \texttt{Frustration: 9} \\
\texttt{Apathy: 0} \\
\texttt{Fear: 0} \\
\texttt{Agreement: 0} \\
            }
        }
    }
    & & \\ \\

    \theutterance \stepcounter{utterance}  
    & & & \multicolumn{2}{p{0.3\linewidth}}{
        \cellcolor[rgb]{0.9,0.9,0.9}{
            \makecell[{{p{\linewidth}}}]{
                \texttt{\tiny{[GM$|$GM]}}
                \texttt{Frustration: 8} \\
\texttt{Apathy: 0} \\
\texttt{Fear: 0} \\
\texttt{Agreement: 0} \\
            }
        }
    }
    & & \\ \\

    \theutterance \stepcounter{utterance}  
    & & & \multicolumn{2}{p{0.3\linewidth}}{
        \cellcolor[rgb]{0.9,0.9,0.9}{
            \makecell[{{p{\linewidth}}}]{
                \texttt{\tiny{[GM$|$GM]}}
                \texttt{game\_result = LOSE} \\
            }
        }
    }
    & & \\ \\

    \theutterance \stepcounter{utterance}  
    & & \multicolumn{4}{p{0.6\linewidth}}{
        \cellcolor[rgb]{0.95,0.95,0.95}{
            \makecell[{{p{\linewidth}}}]{
                \texttt{\tiny{[P1$\leftrightarrow$P]}}
                \texttt{Post{-}episode Reflexion (not sent to GM):} \\
\texttt{New plan: My previous attempt failed because I assigned a non{-}zero score to "Fear" (2), which was not the most accurate interpretation of Lorelei's emotional state given the context of a heated political debate where she is actively arguing her points rather than expressing terror. While "Frustration" was correctly identified as the primary emotion, the inclusion of "Fear" diluted the accuracy of the prediction. Next time, I will strictly analyze the dialogue for the *dominant* emotional tone and assign a score of 0 to any emotion that is not clearly supported by the text, even if it seems remotely plausible. Specifically, I will recognize that Lorelei's tone is one of passionate disagreement and exasperation, not fear, and ensure only "Frustration" receives a high score while "Apathy," "Fear," and "Agreement" are all scored as 0.} \\
            }
        }
    }
    & \\ \\

\end{supertabular}
}

\end{document}
